% ================================================================
% SECTION 1 — INTRODUCTION
% Target: ~4 pages
% ================================================================

\section{Introduction}
\label{sec:introduction}

\subsection{Problem and Motivation}
\label{sec:prob}

Freight transport is the circulatory system of any modern economy.
In South Africa, where the geography of production is separated from
the geography of export by hundreds of kilometres of demanding
terrain, the efficiency of freight logistics determines the
competitiveness of entire industries~\parencite{worldbank2023lpi}.
Transnet Freight Rail (TFR) operates roughly 20\,000 kilometres of
rail network and, at its historical peak, moved more than 200 million
tonnes of freight per year across six major export and general freight
corridors~\parencite{havenga2021}.
The National Corridor (NATCOR) is the country's principal general
freight artery, running from Johannesburg and Pretoria through the
KwaZulu-Natal midlands to the ports of Durban and Richards Bay ---
the busiest container and dry-bulk export gateways on the African
continent~\parencite{dpme2022}.
In a well-functioning logistics system, NATCOR absorbs the dominant
share of containerised imports, manufactured exports, and agricultural
commodities that flow between the interior and the sea~\parencite{havenga2021}.

South Africa's rail freight system has, however, been under
structural stress for decades.
Since the deregulation of road freight in the late 1980s, market
share has shifted persistently from rail to road, and the
infrastructure and rolling stock of TFR have deteriorated under a
combination of underinvestment, load-shedding, cable theft, and
operational inefficiency~\parencite{stander2005, havenga2021}.
By 2021, rail's share of the general freight market had fallen to
levels not seen since the early post-deregulation period, and the
country's rail network was described in policy literature as
approaching a ``tipping point of collapse'' from which recovery would
be extremely costly~\parencite{havenga2021}.
The Presidential State Owned Entities Council noted in 2022 that
Transnet's operational performance had deteriorated significantly,
with locomotive availability falling below 50\% in some quarters
and cable theft incidents rising sharply~\parencite{transnet2023ar}.
Against this backdrop of fragility, the system's capacity to absorb
large exogenous shocks became critically important.

In April 2022, KwaZulu-Natal (KZN) was struck by the most destructive
flooding event in the province's recorded history.
More than 400 people were killed, over 40\,000 displaced, and
infrastructure damage was estimated at tens of billions of
rands~\parencite{ndmc2022}.
The Durban port --- South Africa's largest container terminal ---
was severely disrupted for weeks, with the South African Association
of Freight Forwarders reporting container backlogs of several weeks
in the immediate aftermath~\parencite{saaff2022}.
Rail embankments were washed away, signalling infrastructure was
destroyed, and sections of the NATCOR mainline were rendered
impassable~\parencite{transnet2022floods}.
The shock was sudden, geographically concentrated, and affected a
system already operating under chronic strain.

The question this thesis seeks to answer is deceptively simple: by
how much did the 2022 KZN floods reduce freight volumes on the NATCOR
corridor, and for how long did the effect persist?
The challenge lies not in the question itself but in answering it
rigorously.
NATCOR freight volumes move with the business cycle, respond to fuel
prices, load-shedding, security incidents, and port congestion ---
all of which continued to evolve after April 2022~\parencite{havenga2021, transnet2023ar}.
A naive before-and-after comparison would confound the flood's causal
impact with all of these concurrent trends~\parencite{abadie2010}.
What is needed is a credible counterfactual: an estimate of what
NATCOR volumes \emph{would have been} in the absence of the flood.

\subsection{Why Time-Series Prediction Alone Is Not Enough}
\label{sec:why_causal}

The standard toolkit for freight demand analysis consists of
time-series forecasting models --- autoregressive models (ARIMA and
its variants), exponential smoothing, and increasingly, machine
learning approaches such as gradient boosting and recurrent neural
networks~\parencite{box2015time, hyndman2018forecasting}.
These methods are well suited to predicting the next value in a
series, but they are fundamentally associative.
They learn statistical patterns from historical data without encoding
any notion of causation, and they are particularly ill-equipped to
handle structural breaks --- sudden, permanent changes in the
data-generating process caused by an intervention or
shock~\parencite{pearl2009}.

When a forecasting model is trained on pre-flood data and then asked
to produce a counterfactual for the post-flood period, it produces
a prediction that reflects the average historical behaviour of the
series.
It cannot know whether the deviation between its forecast and the
actual outcome is caused by the flood, by a simultaneous deterioration
in locomotive availability, by a broader economic slowdown, or by
some combination of all three.
The model conflates correlation with causation~\parencite{pearl2009}.
A decision-maker relying on such an analysis might attribute to the
flood a decline that was already underway, or conversely, dismiss a
genuine flood effect because other factors happened to push volumes
upward at the same time~\parencite{abadie2010}.

The limitations of associative prediction become even more pronounced
in the presence of temporal dynamics that are characteristic of
freight systems.
Delayed effects are common: infrastructure repairs take months,
shipper behaviour adjusts with a lag, and diverted freight flows
may not return to rail even after physical repairs are
complete~\parencite{cavallo2013}.
Seasonality interacts with a shock in non-trivial ways: a flood that
strikes during a peak season will produce a larger short-run volume
loss than the same flood in a trough month, but the long-run effect
on shipper confidence may be symmetric~\parencite{hyndman2018forecasting}.
Structural changes --- such as the loss of a key rail customer or
the permanent diversion of freight to road hauliers --- can make
pre-flood statistical relationships obsolete as a guide to
post-flood counterfactual behaviour~\parencite{stander2005}.
None of these dynamics can be resolved by forecasting alone.

What decision-makers need is not a better prediction of what
\emph{will} happen, but a credible estimate of what \emph{would have
happened} under a different policy or in the absence of a
shock~\parencite{pearl2009, abadie2010}.
This is the province of causal inference.

\subsection{Research Gap}
\label{sec:gap}

Despite the strategic importance of South Africa's rail freight
corridors, empirical studies that move beyond associative
description are scarce.
Most existing analyses of South African freight rail either describe
aggregate trends in modal share~\parencite{stander2005, havenga2021} or
develop demand forecasts using regression and time-series models
without a causal identification strategy.
Several important gaps are apparent.

First, prior studies forecast well but remain correlation-based.
They identify co-movement between freight volumes and macroeconomic
or operational variables --- GDP, fuel prices, locomotive
availability --- but they do not establish that these variables
\emph{cause} changes in freight volumes in the interventionist sense
required for policy analysis~\parencite{pearl2009, granger1969}.

Second, existing work tends to use exogenous variables as regressors
without a causal interpretation framework.
Including fuel price as a covariate in a regression is not the same
as estimating the causal effect of a fuel price change on freight
volume; the former conflates the causal effect with confounding
pathways that the regression does not control
for~\parencite{pearl2009}.

Third, and most directly relevant to this thesis, there is no
published causal evaluation of the 2022 KZN flood shock on NATCOR
freight volumes using a counterfactual framework.
The disaster has been described in logistics and insurance
literature~\parencite{ndmc2022, saaff2022}, but the question of how
much volume was actually lost --- relative to a credible benchmark
of what would have occurred without the flood --- has not been
answered with the rigour that causal inference methods now make
possible~\parencite{abadie2003, abadie2010}.

Fourth, robustness across time periods and regime changes has rarely
been assessed.
Freight systems are subject to structural shifts (deregulation,
policy changes, COVID-19 disruptions in 2020--2021) that can make
pre-period statistical relationships unreliable guides to post-period
counterfactual behaviour~\parencite{havenga2021}.
A credible causal analysis must demonstrate that its estimates are
stable across reasonable alternative specifications of the pre-period
and the donor pool~\parencite{benmichael2021}.

\subsection{Research Objective and Questions}
\label{sec:rq}

The overarching objective of this thesis is to develop and apply a
Synthetic Control Method (SCM) framework to estimate the causal
effect of the April 2022 KwaZulu-Natal floods on monthly rail freight
volume on the NATCOR corridor, and to assess the robustness of that
estimate across five distinct causal estimators.

This objective is decomposed into four sub-objectives:
\begin{enumerate}
  \item To construct a clean, longitudinal panel dataset covering six
        South African rail freight corridors at monthly frequency
        over the period January 2015 to December 2024.
  \item To evaluate the performance of five SCM-family estimators
        (Base SCM, ASCM, EN-SCM, MC-SCM, and SDID) in terms of
        pre-period fit quality, causal effect magnitude, and
        robustness to perturbation.
  \item To demonstrate empirically that the adequacy of the donor
        pool is a precondition for the coherence of SCM-based
        estimates, by comparing results obtained with one donor
        corridor against results obtained with five.
  \item To quantify the sustained economic cost of the KZN flood
        shock to NATCOR freight throughput and to assess the
        policy implications for infrastructure resilience planning.
\end{enumerate}

These sub-objectives motivate the following research questions:

\medskip
\noindent\textbf{RQ1:} What is the causal effect of the April 2022
KZN floods on monthly NATCOR rail freight volume, and how does this
estimate vary across five SCM-family estimators?

\medskip
\noindent\textbf{RQ2:} How sensitive are SCM estimates to the size
and composition of the donor pool, and what is the minimum donor pool
required for consistent counterfactual estimation?

\medskip
\noindent\textbf{RQ3:} How robust are the causal estimates to
perturbations in the pre-period window, donor pool composition, and
model hyperparameters?

\medskip
\noindent\textbf{RQ4:} What is the cumulative freight volume loss
attributable to the flood shock, and what does this imply for the
economic cost of infrastructure disruptions on the NATCOR corridor?

\subsection{Contributions and Thesis Structure}
\label{sec:contrib}

This thesis makes three contributions to the literature on causal
inference in transport economics.

\textbf{Methodological:} It applies and systematically compares five
members of the SCM family --- a set of methods that are well
established in economics and political science~\parencite{abadie2003,
abadie2010, benmichael2021, athey2021, arkhangelsky2021} but have
seen limited application in logistics and transport research.
By reporting all five estimators on the same dataset under the same
evaluation protocol, the thesis provides a replicable benchmark for
future SCM-based studies in freight and infrastructure analysis.

\textbf{Empirical:} It provides the first causal estimate of the 2022
KZN flood effect on NATCOR freight volumes using a counterfactual
identification strategy, finding a cumulative volume loss of
approximately 5.1--5.2 million tonnes over the 33 months following
the flood --- a result that is consistent across all five estimators
and survives rigorous placebo and leave-one-out robustness
tests~\parencite{abadie2010}.

\textbf{Diagnostic:} It provides an empirical demonstration of donor
pool inadequacy as a source of estimator divergence, consistent with
warnings in the SCM literature about small donor
pools~\parencite{abadie2003, benmichael2021}.
When a single donor corridor (CAPE) is used, the five methods produce
estimates ranging from $-11.1$\,MT to $+3.7$\,MT --- a span of
nearly 15\,MT that renders any single estimate meaningless.
When a proper five-corridor donor pool is constructed, the estimates
converge to a range of just 0.15\,MT, providing evidence that
donor pool adequacy is a precondition for, rather than a product of,
causal identification in the SCM framework.

The remainder of this thesis is structured as follows.
Section~\ref{sec:literature} reviews the relevant literature on causal
AI for time-series data, the SCM family of estimators, and the South
African freight rail context, positioning this study relative to prior
work.
Section~\ref{sec:methodology} describes the full analytical framework,
including data sources, preprocessing, the five estimators, evaluation
metrics, and robustness procedures.
Section~\ref{sec:results} presents the empirical results, beginning
with the inadequate single-donor analysis and progressing to the
five-donor causal estimates and their robustness.
Section~\ref{sec:discussion} interprets the findings in light of
the research questions, relates them to prior literature, and
discusses the practical implications for infrastructure resilience
policy.
Section~\ref{sec:conclusion} concludes with a summary of
contributions and directions for future research.
